\documentclass[11pt,a4paper]{article}

\usepackage[T1]{fontenc}
\usepackage{lmodern}
\usepackage[margin=23mm]{geometry}
\usepackage{amsmath,amssymb}
\usepackage{booktabs,tabularx}
\usepackage[hidelinks]{hyperref}
\hypersetup{pdftitle={Gauging and ungauging on two vertices}}

\setlength{\parindent}{0pt}
\setlength{\parskip}{0.5em}

\begin{document}
{\LARGE\bfseries Gauging and ungauging on two vertices\par}

A graph with two vertices and one edge is enough to see the difference between
the round-trip and persistent circuits. Both measure a data operator \(U\).
The round trip ends by reading the edge and using its bit to correct the data.
The persistent circuit performs that correction coherently at entry and uses
the resulting coordinates throughout verification.

\section*{Common preparation}

The two vertex ancillas are \(a_0,a_1\), and the edge ancilla is \(A\). Vertex
\(0\) is the root. The data register is separate from the graph and starts in
an encoded state \(|\psi\rangle\). The operators \(U_0\) and \(U_1\) act on
disjoint groups of data qubits, and their product \(U\) preserves the code
space. They satisfy

\[
U_0^2=U_1^2=I,\qquad U_0U_1=U_1U_0=U.
\]

Prepare the vertices in \(|+\rangle\) and the edge in \(|0\rangle\). A CNOT
from each vertex into \(A\) writes the XOR of the two vertex bits on the edge.
Then use \(a_0\) to control \(U_0\) and \(a_1\) to control \(U_1\). These
two controlled operations are the enrichment \(\Omega\). The joint state is

\begin{equation}
\label{eq:enriched-state}
\frac12\Big[
 |00\rangle_a|0\rangle_A|\psi\rangle
+|01\rangle_a|1\rangle_A U_1|\psi\rangle
+|10\rangle_a|1\rangle_A U_0|\psi\rangle
+|11\rangle_a|0\rangle_A U|\psi\rangle
\Big].
\end{equation}

In both circuits, the vertices are eventually measured in \(X\). Record
their bits \(b_0,b_1\) and write \(x_j=(-1)^{b_j}\). The parity of the
two reports gives the sign \(r=x_0x_1\). For a vertex term labelled by
a bit \(t\), the measurement contributes the amplitude
\(\langle b_j|H|t\rangle=x_j^t/\sqrt2\). Thus the terms \(00,01,10,11\)
in \eqref{eq:enriched-state} acquire factors \(1,x_1,x_0,r\), respectively, and a common factor \(1/2\).

Write \(P_r=(I+rU)/2\). Whenever \(P_r|\psi\rangle\) is nonzero, denote its
normalized state by \(|\psi_r\rangle\). It obeys
\(U|\psi_r\rangle=r|\psi_r\rangle\). In the clean positive branch used for
state preparation, \(r=+1\).

\section*{Round trip}

\subsection*{1. Measure the vertices}

Measure \(a_0,a_1\) in \(X\) directly after \eqref{eq:enriched-state}. The terms \(00,11\) both
have edge label \(0\), while \(01,10\) both have edge label \(1\). Adding
the contributions for each edge gives the unnormalized state

\[
\frac14\Big[
 |0\rangle_A (I+rU)|\psi\rangle
+|1\rangle_A (x_1U_1+x_0U_0)|\psi\rangle
\Big].
\]

On the edge-\(1\) term, \(x_1U_1+x_0U_0=x_1U_1(I+rU)\), since
\(x_0=x_1r\) and \(U_1U=U_0\). The known sign \(x_1\) on that term can be
removed by \(Z_A^{b_1}\). With the sign set this way, the normalized state
in the \textbf{gauged coordinates} is

\begin{equation}
\label{eq:gauged-state}
\frac{|0\rangle_A|\psi_r\rangle+|1\rangle_A U_1|\psi_r\rangle}{\sqrt2}.
\end{equation}

\subsection*{2. Apply \(V\), then measure the ordinary data checks and \(X_A\)}

The coordinate change is \(V=C_A(U_1)\), with the \textbf{edge qubit} as the
control. On the edge-\(1\) term in \eqref{eq:gauged-state}, it applies another \(U_1\).
The two applications cancel, so the state in the \textbf{corrected coordinates} is

\begin{equation}
\label{eq:corrected-state}
V\frac{|0\rangle_A|\psi_r\rangle+|1\rangle_A U_1|\psi_r\rangle}{\sqrt2}
=|+\rangle_A|\psi_r\rangle.
\end{equation}

Now measure the ordinary data stabilizers. They check whether the data
lies in the original code space. For an ordinary stabilizer \(S\), the
clean outcome is \(+1\) because \(S|\psi_r\rangle=|\psi_r\rangle\).
A separate check of \(U\) is needed to tell its positive and negative
sectors apart.

Then measure \(X_A\) \textbf{using a helper}, so the edge stays available for
the next gate. Prepare the helper in \(|+\rangle\), apply a CNOT from
the helper to \(A\), then read the helper in \(X\). Here the edge itself is
in \(|+\rangle_A\). Since \(X_A|+\rangle_A=|+\rangle_A\), the CNOT leaves
both qubits in \(|+\rangle\) and the helper reports bit \(0\), meaning
\(X_A=+1\). The state after these clean checks is still \eqref{eq:corrected-state}.

\subsection*{3. Apply \(V\) again, then measure the root Gauss check}

The second \(V\) takes \eqref{eq:corrected-state} back to \eqref{eq:gauged-state}. This return lets the circuit use the
root Gauss check \(G_0=U_0X_A\), whose data part acts on just the root group.
The check at the other vertex is \(G_1=U_1X_A\). On \eqref{eq:gauged-state}, \(G_1\) swaps the
two terms, so its eigenvalue is \(+1\). This is the effect of the \(X_A=+1\)
check in the corrected coordinates.

Since \(G_0G_1=U_0U_1=U\), the root result gives the sign of \(U\) when
\(G_1=+1\). The same result follows directly by applying \(G_0\) to \eqref{eq:gauged-state}:

\[
G_0
\frac{|0\rangle_A|\psi_r\rangle+|1\rangle_AU_1|\psi_r\rangle}{\sqrt2}
=
\frac{|1\rangle_A U_0|\psi_r\rangle+|0\rangle_AU|\psi_r\rangle}{\sqrt2}
=
r\,\frac{|0\rangle_A|\psi_r\rangle+|1\rangle_AU_1|\psi_r\rangle}{\sqrt2}.
\]

The last equality uses \(U_0|\psi_r\rangle=U_1U|\psi_r\rangle
=rU_1|\psi_r\rangle\). The root check reports \(r\), and the clean state
is still \eqref{eq:gauged-state}. Each round of verification follows the same pair of
coordinate changes.

\subsection*{4. Ungauge by reading the edge in \(Z\)}

At the end of verification, measure the edge itself in \(Z\). The result
\(z=0\) selects \(|\psi_r\rangle\). The result \(z=1\) selects
\(U_1|\psi_r\rangle\). Apply the data correction \(U_1^z\), which gives
\(U_1^zU_1^z|\psi_r\rangle=|\psi_r\rangle\) in either case.

The \(X_A\) check during verification reads the edge's \(X\) eigenvalue
through a helper. This final \(Z\) readout consumes the edge and tells the
controller which data correction to apply.

\section*{Persistent coordinates}

\subsection*{1. Apply \(V\) before measuring the vertices}

Start from \eqref{eq:enriched-state}. This time apply \(V=C_A(U_1)\) while the vertices are still
unmeasured. It changes the data operation on the \(01\) term from \(U_1\)
to \(I\), and on the \(10\) term from \(U_0\) to \(U_1U_0=U\). Both of
these terms have edge label \(1\). The terms with edge label \(0\) retain
their data operations. Thus the state becomes

\[
\frac12\Big[
 |00\rangle_a|0\rangle_A|\psi\rangle
+|01\rangle_a|1\rangle_A|\psi\rangle
+|10\rangle_a|1\rangle_A U|\psi\rangle
+|11\rangle_a|0\rangle_A U|\psi\rangle
\Big].
\]

Now measure the vertices in \(X\), using the same bit-to-sign rule.
Grouping by edge label as before gives

\[
\begin{aligned}
&\frac14\Big[
 |0\rangle_A (I+rU)|\psi\rangle
+|1\rangle_A (x_1 I+x_0U)|\psi\rangle
\Big]\\
&\qquad=
\frac14\bigl(|0\rangle_A+x_1|1\rangle_A\bigr)
\otimes(I+rU)|\psi\rangle .
\end{aligned}
\]

The factorization uses \(x_0=x_1r\). Setting the known edge sign with
\(Z_A^{b_1}\) and normalizing gives \(|+\rangle_A|\psi_r\rangle\), the
same product state as \eqref{eq:corrected-state}. Verification starts in the corrected
coordinates.

\subsection*{2. Measure the ordinary data checks, \(X_A\), and the transformed root check}

Measure \(S\) on the data and \(X_A\) through the helper, just as in the
round trip. Both report \(+1\) in this clean example. For the root
measurement, stay in the corrected coordinates and measure
\(VG_0V^\dagger\).

To see what this operator does, take an edge basis term \(|z\rangle_A\).
The right-hand \(V^\dagger\) contributes \(U_1^z\). Then \(G_0=U_0X_A\)
applies \(U_0\) and flips the edge bit to \(z\oplus1\). The left-hand
\(V\) contributes \(U_1^{z\oplus1}\). The combined data operation is

\[
U_1^{z\oplus1}U_0U_1^z=U_0U_1=U.
\]

Exactly one of the two edge controls contributes \(U_1\), because the
edge flips between them. Therefore

\[
VG_0V^\dagger=UX_A,\qquad
(UX_A)|+\rangle_A|\psi_r\rangle=r\,|+\rangle_A|\psi_r\rangle.
\]

The edge contributes \(+1\), so the measurement reads the data sign
\(r\). The root check now acts on \textbf{both} data groups and the edge.
Subsequent verification rounds use these same corrected coordinates.

\subsection*{3. Discard the edge}

The final state is \(|+\rangle_A|\psi_r\rangle\). Its edge-\(0\) and
edge-\(1\) terms both carry the same data state. The stock persistent
circuit measures the edge in \(Z\) for disposal. Either outcome leaves
\(|\psi_r\rangle\) on the data, so the edge report can be ignored.

\section*{Measurement order}

\begin{tabularx}{\textwidth}{@{}l>{\raggedright\arraybackslash}X>{\raggedright\arraybackslash}X@{}}
\toprule
\textbf{Stage} & \textbf{Round trip} & \textbf{Persistent} \\
\midrule
After enrichment \(\Omega\) & Read vertices in \(X\). & Apply \(V\), then read vertices in \(X\). \\
Data and edge checks & Apply \(V\), then measure \(S\) and \(X_A\). & Measure \(S\) and \(X_A\) in the current coordinates. \\
Root check & Apply \(V\) again, then measure \(U_0X_A\). & Measure \(U_0U_1X_A\) in the current coordinates. \\
After verification & Read edge in \(Z\) and apply \(U_1^z\). & Read edge in \(Z\) for disposal. \\
\bottomrule
\end{tabularx}

The full block and its verification schedules are described in
\href{run:Understandable_Inter_CZ.md}{Understanding the internal logical CZ measurement}.
The two-vertex example follows the ideal measurement sequence and the
data byproduct through both circuits.

\end{document}
